% Sections 14 to 16, plus references.

\section{Threat Model: The Honest Register}\label{sec:threats}

A whitepaper that does not list its risks is a sales document; this one
is not. The threats are ranked by severity, with their mitigations and,
because a mitigated risk is not a vanished risk, their residual status
stated in plain words. That discipline is not decorative: the previous
sections showed that a social specification judged solid could fall
before its own simulation; the present section applies the same
requirement to the whole assembly, before an external audit takes it
over again, phase by phase, on the roadmap of Section~\ref{sec:roadmap}.

\begin{table}[htbp]
\centering
\small
\caption{Threat register: mitigations and residual status.}
\label{tab:threats}
\begin{tabularx}{\textwidth}{@{}L{3.4cm}L{1.5cm}YY@{}}
\toprule
Threat & Severity & Mitigation & Residual status \\
\midrule
mining reorganization (majority attack) & high & Ring checkpoints: everything signed is immutable; the attack only prospers on the six-second window & reduced; double spending remains possible on the unsigned window \\
Ring collusion (55 seats) & high & weighted draw over a large population, rotation each era, term limit of three eras, total stripping on the first signed fork & low in steady regime; real during the first eras \\
massive reputation purchase & high & the Tenure layer is not monetizable; Echo is capped and weighted; concentration index monitored; rules R1 to R4 & structurally blocked by time; the assumed wall \\
false humanities (sybil on Embers) & medium & sponsorship capped at two per year, sponsor Kleos at stake, presence proof each era & depends on the size of the established web; the first years are the most exposed \\
DAG on a CryptoNote base (engineering) & high & the seams specified and implemented, $k = 8$, order resolution, decoy sampling, pruning; isolated prototype, dedicated audit, capacity budget with measured exit criteria & difficulty, not uncertainty; the long pole of the roadmap \\
regulatory frame (MiCA, Travel Rule) & high & the Lumen: provable disclosure, native data minimization; payment-currency strategy & structural, not eliminable; readability is the maximum reachable \\
strict-cap failure (long-term security) & medium & payment volume as the fee resource; public security budget tracking; constitutional tail safety valve & the assumed wager of the symbolic contract \\
the false thousand (Corona sybil) & medium & bond per senate seat; hardware attestation, one chip one node; concentration caps & bounded, auditable, staged activation \\
false diversity (one base model) & medium & model-family caps; random-draw inspection; three blind implementations in the Arena & a standing watch of governance, not a solved problem \\
machine cost (honest operators priced out) & medium & few well-paid seats; bounded treasury budgets; sunset forces a re-vote & the wager that useful intelligence pays, tested small before large \\
biometric debt & none & no body data collected, anywhere, ever & none; nothing to leak \\
\bottomrule
\end{tabularx}
\end{table}

Two threats deserve development. Ring collusion first: it is the price
of the finality assembly, and its remedy is not technical but
demographic. The larger the high-reputation population, the more buying a
majority of seats requires buying years of irreproachability from people
who have no reason to sell: the network must therefore cultivate its
reputation class as a common good, from the first eras onward; that is a
policy, not a function. The regulatory risk next: the worst strategy
would be discretion, hiding attracts suspicion, being readable invites
expertise. The documentation of the Lumen, the compliance proofs and the
data-minimization argument are designed to be shown, not endured: the
transparency of the mechanism is the best defense of the opacity of the
data.

\section{The Zero-Defect Method}\label{sec:method}

The question that closes every conversation about this project is the
right one: will you be able to code all of that without a single fault?
The honest answer begins with an admission: nobody can promise it, and
whoever promises it is selling something. Bitcoin itself, the most
scrutinized implementation in the world, nearly failed twice. An integer
overflow created one hundred eighty-four billion bitcoins in one
transaction in 2010, spotted within hours, chain replayed by hand; the
2018 vulnerability would have allowed minting infinite bitcoins, found
by a volunteer auditor, cold, on code examined for ten years. The
CryptoNote base had its own episode in 2017: a flaw in key images
allowed double spending on every chain of the family. The right question
is therefore not whether the fault is possible, it always is; it is
which method catches faults before attackers do.

\subsection{Two captured specimens}

Before the specimens, one event worth recording, because it is the
method working on the method: the first external audit of version 1.1
(September 2026, an independent reviewer) returned twelve findings,
factual and design-level, from a wrong RandomX memory figure to the
unspecified bootstrap of the Ring. Version 1.2 closes the twelve: the
numbers were corrected, the missing mechanisms were specified (the
Genesis Ring ratchet, the identity firewall, the precise revocation of
a Cipher, the capacity budget), the overclaims were cut (the ring of
sixteen now equals Monero's, it does not triple anything; the golden
ratio is now named as an identity, not an argument), and the simulation
now sweeps twelve seeds. The register of what an outside eye catches
that the author's eye does not is the whole argument of layer five
below, and this paper has now been through it once for real.

First specimen, lived on XelisVault on September 5, 2026. The paper
wallet generated a twenty-five word mnemonic seed; imported into the
official NERVA wallet, it restored an address different from the printed
one. The code looked right, re-read twice, and that is precisely the
trap: the bug lived in a reading convention, invisible to re-reading.
The C++ reference groups the seed into four-byte words and writes them
in little-endian order; our TypeScript mirror read them in the opposite
order. The fix was made the way it must always be made: port the
reference line by line, generate forty-eight test vectors, require bit
to bit identity of the two implementations, then re-run the audit. The
lesson exceeds the case: an encoding fault is not found, it is crossed;
only a second implementation, executed in parallel on the same inputs,
reveals it mechanically.

Second specimen, deeper, since it concerns the design itself and not one
line of code: the Kleos simulation attack, told in
Section~\ref{sec:kleos}. A farm of fake profiles crossed the candidacy
threshold for the Ring before the honest network and captured the
fifty-five seats; four corrective rules closed the window, and the same
attack replayed takes no seat, on any of the twelve seeds. The decisive
point is not that the flaw was found: it is that it was found by a
three-hundred-line simulation, before any chain code, at a cost in
hours of computation and not in years of reputation. The deterministic
simulation of a specification is the cheapest machine there is for
catching design faults; every new rule of the protocol will be treated
that way before implementation, with numeric invariants, a seed sweep
and a regression run over three attack sets.

\subsection{The six-layer verification protocol}

Those two specimens define the method better than a treatise would; it
fits in six layers, each extinguishing an entire class of faults rather
than one particular bug. A bug of an extinguished class does not come
back; a fixed bug always comes back. First layer: cross double
implementation on test vectors, which reveals encoding and convention
faults. Second layer: deterministic simulation of the social and
economic rules before any code, which reveals design flaws. Third layer:
reproducible builds and continuous massive random testing, which reveal
non-determinisms and unexpected states. Fourth layer: systematic
cross-review between the project holder and his assistant, two pairs of
eyes, two reading conventions. Fifth layer: external audit of the
differential at every phase, because a fresh auditor does not have the
author's blind spots. Sixth layer: measurable GO and NO-GO exit criteria
at every roadmap phase, with the rule that a NO-GO remains an acceptable,
published result.

The implementation discipline adds to the six layers: never a main
network on release-candidate foundations; multi-platform distribution
from day one, five signed targets; refusal of unaudited exotic
cryptographic primitives, stated since the first version of this
specification and held since; publication of all code from the
specification phase onward, because hidden code is code nobody has
looked at. The method does not promise perfection; it promises the
shortest possible reaction time between the appearance of a fault and
its capture, and that is the only promise an honest engineer can make.

\section{Roadmap and Conclusion}\label{sec:roadmap}

The first version of this paper pinned an eighteen-month calendar to
genesis, and the external audit called it what it was: a schedule for a
skeleton of the scope. Monero, Kaspa and Xelis each spent years on one
of this project's subsystems. The honest form of a roadmap is
criteria-driven: the phases below carry month ranges as planning aids
for a team of at least three full-time engineers, but the binding term
is the exit criterion of each phase, and genesis happens when every
criterion is green, not when a Gantt chart says so. The estimate the
paper stands behind is twenty-four months of full-time work to the
earliest responsible genesis, with the DAG-plus-rings prototype as the
long pole, and a published NO-GO as an acceptable outcome at every
step. The base is an assembly of proven foundations, the CryptoNote
lineage for the private sphere and the public literature of
fast-converging DAGs for the ordering layer; everything that is new,
Kleos, the Ring, the Lumen, Cipher perimeters, Mandates, the Corona, is
deterministic, testable, auditable code, with no exotic primitive
before external audit. Every phase carries one measurable deliverable
and one binary exit criterion; an unmet criterion is not a shameful
failure, it is steering data, and the phase replays until it is met.

\begin{table}[htbp]
\centering
\small
\caption{Twenty-four months, six phases, one binary exit criterion at
every step; genesis when green, not when due.}
\label{tab:roadmap}
\begin{tabularx}{\textwidth}{@{}L{2.5cm}L{1.4cm}YY@{}}
\toprule
Phase & Months & Measurable deliverables & Exit criterion \\
\midrule
1, specification & M1-M3 & written protocol, numbered architecture decisions, format sketches, extended social-core simulation with the seed sweep, red-team bounties of Corona stage 0 & specification read by two external readers; first audit findings closed \\
2, DAG prototype & M4-M8 & dev network mining a CPU DAG at two seconds, stem propagation, the capacity budget measured & twenty-four hours without unexpected reorganization; 100 tx/s sustained on the reference node \\
3, Ring and Kleos v0 & M9-M13 & signed checkpoints, Deed score computed, era rotation replayed, the bootstrap ratchet drilled & finality under six seconds over one hundred thousand replayed blocks \\
4, identities and Lumen & M14-M18 & Embers with sponsorship, Ciphers with warrants and fallback Mandates, bounded viewing keys, external audit of the differential & audit without a critical flaw \\
5, public test network & M19-M23 & faucet, explorer, one hundred nodes, agent load simulator, upgrade ritual repeated three times, Arena drills & three consecutive upgrades without incident; 300 tx/s burst validated or the target dropped \\
6, genesis & M24 at the earliest & public ceremony, signed executables for five platforms, Lumen documentation for authorities published & all previous criteria, without exception \\
\bottomrule
\end{tabularx}
\end{table}

Three ecosystem commitments accompany the calendar. The code is public
from the specification phase onward, and this whitepaper lives next to
it in the same repository, revised like code. The XelisVault ecosystem,
register, tickets, paper wallet, becomes multi-chain at phase five: a
feature flag, not a rewrite, NERVA first, ANTUMBRA after. Finally, the
project identities, domain, repositories, social networks, are reserved
before phase two, so that nothing lends itself to confusion on genesis
day. The roadmap reads as a promise of operational humility:
criteria-driven work, every step verifiable by a third party, none of it
depending on a miracle of cryptography, and a calendar that bends to the
evidence instead of the evidence bending to the calendar.

A faithfully kept shadow, a ring of light anyone can verify on demand:
privacy as a right, proof as a courtesy, and time as the only justice
of the peace that money cannot corrupt.

Such is the public specification of the ANTUMBRA network. The demand of
2026 is real and measurable: agents that pay, bots to distinguish from
persons, authorities to whom one shows proofs rather than promises,
users who want the instant. The answer fits one image, the one that
gives the network its name: the antumbra, the zone of the annular
eclipse where a source too small to be hidden lets a ring of light show
around the shadow. The core of the network is that shadow, balances,
identities, flows, kept by construction; around it, the ring of
verification, reputation, finality, compliance proofs, which anyone can
switch on without exposing anyone other than oneself. Around the ring
itself, one more loop is sketched and deliberately left unlit until its
stage criteria are met: the Corona, machine governance under human
veto, the network learning to prune its own orchard while the humans
keep the fruit. The specification was written to be attacked before
being implemented, and that is the only honest way to write a
whitepaper. To work.

\section*{References}
\addcontentsline{toc}{section}{References}
\begin{enumerate}[leftmargin=1.6em, itemsep=1.5pt, topsep=4pt]
\item S. Nakamoto. \emph{Bitcoin: A Peer-to-Peer Electronic Cash System}. bitcoin.org, 2008.
\item N. van Saberhagen. \emph{CryptoNote v 2.0}. 2013.
\item tevador et al. \emph{RandomX Proof of Work}. randomx.io, 2019.
\item Y. Sompolinsky, S. Wyborski, Y. Zohar. \emph{PHANTOM and GHOSTDAG: A Scalable Generalization of Nakamoto Consensus}. 2021.
\item B. B\"unz, S. Agrawal, S. Boneh, M. Weston. \emph{Bulletproofs: Short Proofs for Confidential Transactions and More}. IEEE S\&P, 2018.
\item R. Can B\"unz, S. Agrawal, S. Boneh, M. Weston et al. \emph{Zero to Bulletproofs+}. 2020.
\item C. Komlo, I. Goldberg. \emph{FROST: Flexible Round-Optimized Schnorr Threshold Signatures}. 2020.
\item S. Noether, B. Goodell et al. \emph{CLSAG: Compressed Linkable Spontaneous Anonymous Group Signatures}. 2020.
\item S. B. Venkatakrishnan, G. Fanti, M. Zhou, K. Cai, P. Viswanath et al. \emph{Dandelion++: Lightweight Cryptocurrency Networking with Formal Anonymity Guarantees}. ACM ToMPECS, 2018.
\item T. Pedersen. \emph{Non-Interactive and Information-Theoretic Secure Verifiable Secret Sharing}. EUROCRYPT, 1991.
\item M. Castro, B. Liskov. \emph{Practical Byzantine Fault Tolerance}. OSDI, 1999.
\item D. Hopwood et al. \emph{Zcash Protocol Specification}. zips.z.cash, ongoing.
\item Regulation (EU) 2016/679 (General Data Protection Regulation), Article 5(1)(c), data minimization.
\item Regulation (EU) 2023/1114 (Markets in Crypto-Assets, MiCA).
\end{enumerate}
